%---------------------------------------------------------------------
% Part three: hadith. Four questions.
%
% The bab titles the papers name -- حفظ اللسان، الاستیذان، الحب فی اللہ،
% البر والصلۃ -- are the prescribed chapter headings of the course text,
% which could not be obtained. So each hadith below is attributed to the
% collection that actually grades it (Bukhari, Muslim, Tirmidhi), which
% is what an AIOU script is expected to name in any case.
%---------------------------------------------------------------------

\section*{حصہ سوم --- احادیث}
\markboth{احادیث}{}

%---------------------------------------------------------------------
\begin{sawal}{سوال \textenglish{6} \ \textnormal{\small(تین پرچوں میں --- بہار \textenglish{2012} س\,\textenglish{3}، بہار \textenglish{2013} س\,\textenglish{4}، خزاں \textenglish{2021} س\,\textenglish{5})}}
''حفظ اللسان من الغیبۃ والشتم'' --- احادیث کی روشنی میں زبان کی حفاظت پر تفصیلی
نوٹ لکھیے۔
\end{sawal}

\begin{hal}
یہ اس حصے کا \textbf{سب سے زیادہ دہرایا جانے والا} موضوع ہے --- چار میں سے تین
پرچوں میں آیا۔

\medskip
\textbf{قرآنی بنیاد۔} ہر لفظ لکھا جا رہا ہے: \ar{مَا يَلْفِظُ مِن قَوْلٍ إِلَّا
لَدَيْهِ رَقِيبٌ عَتِيدٌ} (ق، \textenglish{18})۔ اور غیبت کی تصویر یوں کھینچی
گئی:

\ayat{وَلَا يَغْتَب بَّعْضُكُم بَعْضًا أَيُحِبُّ أَحَدُكُمْ أَن يَأْكُلَ
لَحْمَ أَخِيهِ مَيْتًا فَكَرِهْتُمُوهُ}{الحجرات، \textenglish{12}}
\tarjuma{اور تم میں سے کوئی کسی کی غیبت نہ کرے۔ کیا تم میں سے کوئی پسند کرے گا
کہ اپنے مرے ہوئے بھائی کا گوشت کھائے؟ تم اسے ناپسند کرو گے۔}

\medskip
\textbf{بنیادی احادیث:}

\hadees{مَنْ كَانَ يُؤْمِنُ بِاللهِ وَالْيَوْمِ الْآخِرِ فَلْيَقُلْ خَيْرًا
أَوْ لِيَصْمُتْ}{بخاری و مسلم}
\tarjuma{جو اللہ اور یومِ آخرت پر ایمان رکھتا ہو، وہ اچھی بات کہے یا خاموش رہے۔}

\smallskip
یعنی گفتگو میں تیسرا راستہ ہے ہی نہیں --- یا بھلائی، یا سکوت۔

\hadees{الْمُسْلِمُ مَنْ سَلِمَ الْمُسْلِمُونَ مِنْ لِسَانِهِ وَيَدِهِ}{بخاری
و مسلم}
\tarjuma{مسلمان وہ ہے جس کی زبان اور ہاتھ سے دوسرے مسلمان محفوظ رہیں۔}

\smallskip
غور کیجیے کہ \emph{زبان کو ہاتھ سے پہلے} رکھا گیا۔

\hadees{سِبَابُ الْمُسْلِمِ فُسُوقٌ وَقِتَالُهُ كُفْرٌ}{بخاری و مسلم ---
حضرت عبد اللہ بن مسعودؓ \ \ [بہار \textenglish{2013}]}
\tarjuma{مسلمان کو گالی دینا فسق ہے اور اس سے لڑنا کفر ہے۔}
\textbf{تشریح:} ''کفر'' سے مراد یہاں دین سے خروج نہیں بلکہ \emph{کفرانِ نعمت}
اور شدید وعید ہے --- یہ نکتہ لکھنا ضروری ہے۔

\medskip
\textbf{غیبت کی تعریف --- خود نبی کریم \saw کی زبانی۔} \ar{ذِكْرُكَ أَخَاكَ
بِمَا يَكْرَهُ} --- اپنے بھائی کا ایسا ذکر جو اسے ناگوار ہو۔ پوچھا گیا: اگر وہ
بات اس میں ہو تو؟ فرمایا: \ar{إِنْ كَانَ فِيهِ مَا تَقُولُ فَقَدِ
اغْتَبْتَهُ، وَإِنْ لَمْ يَكُنْ فِيهِ فَقَدْ بَهَتَّهُ} --- اگر ہو تو غیبت،
نہ ہو تو بہتان (مسلم)۔ \textbf{یہی غیبت اور بہتان کا فرق ہے} --- پوچھا جا چکا
ہے۔

\medskip
\textbf{ضمانت:} \ar{مَنْ يَضْمَنْ لِي مَا بَيْنَ لَحْيَيْهِ وَمَا بَيْنَ
رِجْلَيْهِ أَضْمَنْ لَهُ الْجَنَّةَ} --- زبان اور شرمگاہ کی ضمانت دے تو جنت کی
ضمانت (بخاری)۔

\medskip
\textbf{زبان کے آفات:} غیبت، بہتان، چغلی، جھوٹ، گالی و طعنہ، مذاق اڑانا۔
\textbf{غیبت نہیں:} مظلوم کی فریاد، مسئلہ پوچھنے کو واقعہ بیان کرنا، نقصان سے
خبردار کرنا، اور کھلم کھلا گناہ کرنے والے کا ذکر۔
\end{hal}

%---------------------------------------------------------------------
\begin{sawal}{سوال \textenglish{7} \ \textnormal{\small(تین پرچوں میں --- بہار \textenglish{2012} س\,\textenglish{7}، بہار \textenglish{2013} س\,\textenglish{3}، خزاں \textenglish{2021} س\,\textenglish{3})}}
''باب البر والصلۃ'' کا خلاصہ لکھیے اور اس باب کی احادیث کا ترجمہ و تشریح کیجیے۔
\end{sawal}

\begin{hal}
\textbf{مفہوم۔} \emph{البِرّ} یعنی ہر قسم کی نیکی اور حسنِ سلوک؛ \emph{الصِّلۃ}
یعنی رشتہ جوڑنا، خاص طور پر صلہ رحمی۔ باب کا مرکزی خیال یہ ہے کہ ایمان کا اظہار
عبادت کے ساتھ ساتھ \emph{رشتوں کے برتاؤ} میں ہوتا ہے۔

\medskip
\begin{nukta}[یہ حدیث دو پرچوں میں لفظ بہ لفظ آئی --- اسے ضرور یاد کیجیے]
\hadees{مَنْ أَحَبَّ أَنْ يُبْسَطَ لَهُ فِي رِزْقِهِ وَيُنْسَأَ لَهُ فِي
أَثَرِهِ فَلْيَصِلْ رَحِمَهُ}{بخاری --- حضرت انسؓ \ \ [بہار \textenglish{2013}
اور خزاں \textenglish{2021} دونوں میں]}
\tarjuma{جو شخص چاہتا ہے کہ اس کے رزق میں کشادگی ہو اور اس کی عمر میں اضافہ کیا
جائے، وہ صلہ رحمی کرے۔}

\smallskip
\textbf{تشریح:} ''اثر'' سے مراد عمر یا نشانِ زندگی ہے۔ عمر بڑھنے کا مطلب
علماء نے دو طرح بیان کیا: \textbf{(۱)} عمر میں برکت --- کہ کم وقت میں زیادہ کام
ہو جائے، اور \textbf{(۲)} تقدیرِ معلق کا بدلنا۔ چنانچہ یہ حدیث ''تقدیر مبرم و
معلق'' والے سوال سے براہِ راست جڑی ہوئی ہے۔
\end{nukta}

\medskip
\textbf{اسی باب کی دیگر احادیث:}

\hadees{لَيْسَ الْوَاصِلُ بِالْمُكَافِئِ، وَلَكِنِ الْوَاصِلُ الَّذِي إِذَا
قُطِعَتْ رَحِمُهُ وَصَلَهَا}{بخاری --- حضرت عبد اللہ بن عمروؓ}
\tarjuma{صلہ رحمی کرنے والا وہ نہیں جو بدلہ چکائے، بلکہ وہ ہے جو رشتہ توڑے
جانے پر بھی جوڑے۔}
\textbf{تشریح:} یعنی صلہ رحمی \emph{لین دین} نہیں، \emph{یک طرفہ احسان} ہے۔

\medskip
\textbf{والدین کا مقام۔} پوچھا گیا: میرے حسنِ سلوک کا سب سے زیادہ حق دار کون
ہے؟ فرمایا: ''تیری ماں''، پھر ''تیری ماں''، پھر ''تیری ماں''، پھر ''تیرا باپ''
(بخاری و مسلم)۔ قرآن: \ar{وَبِالْوَالِدَيْنِ إِحْسَانًا}۔

\medskip
\textbf{قطع رحمی کی وعید۔} \ar{لَا يَدْخُلُ الْجَنَّةَ قَاطِعٌ} --- رشتہ توڑنے
والا جنت میں داخل نہ ہو گا (بخاری و مسلم)۔ اور بھائی سے تین دن سے زیادہ قطع
تعلقی حرام ہے (بخاری و مسلم --- خزاں \textenglish{2021})۔

\medskip
\textbf{خلاصہ۔} نیکی کی جامع تعریف بھی اسی باب میں ہے: \ar{الْبِرُّ حُسْنُ
الْخُلُقِ} --- نیکی خوش خلقی ہی کا نام ہے (مسلم)۔
\end{hal}

%---------------------------------------------------------------------
\begin{sawal}{سوال \textenglish{8} \ \textnormal{\small(بہار \textenglish{2013}، سوال \textenglish{4} اور خزاں \textenglish{2021}، سوال \textenglish{1})}}
''الحب فی اللہ'' اور ''الاستیذان'' پر احادیث کی روشنی میں نوٹ لکھیے، نیز
مسلمانوں کے باہمی حقوق بیان کیجیے۔
\end{sawal}

\begin{hal}
\textbf{(الف) الحب فی اللہ۔} یعنی محبت کی بنیاد نہ نسب ہو، نہ مال، نہ مفاد ---
صرف اللہ کی رضا۔

\begin{itemize}\setlength{\itemsep}{2pt}
  \item جن سات آدمیوں کو قیامت کے دن عرش کا سایہ ملے گا، ان میں
        \ar{وَرَجُلَانِ تَحَابَّا فِي اللهِ اجْتَمَعَا عَلَيْهِ وَتَفَرَّقَا
        عَلَيْهِ} بھی ہیں (بخاری و مسلم)۔
  \item \ar{ثَلَاثٌ مَنْ كُنَّ فِيهِ وَجَدَ حَلَاوَةَ الْإِيمَانِ} --- اللہ و
        رسول \saw سب سے محبوب ہوں، بندے سے صرف اللہ کے لیے محبت ہو (بخاری و
        مسلم)۔
  \item \textbf{حدیثِ قدسی:} \ar{وَجَبَتْ مَحَبَّتِي لِلْمُتَحَابِّينَ فِيَّ}
        --- میری محبت ان کے لیے واجب جو میرے لیے آپس میں محبت کرتے ہیں (موطأ
        امام مالک)۔
\end{itemize}

\medskip
\textbf{(ب) الاستیذان --- اجازت لینے کا ادب۔} \ar{لَا تَدْخُلُوا بُيُوتًا
غَيْرَ بُيُوتِكُمْ حَتَّىٰ تَسْتَأْنِسُوا وَتُسَلِّمُوا عَلَىٰ أَهْلِهَا}
(النور، \textenglish{27})۔ حدیث: \ar{الِاسْتِئْذَانُ ثَلَاثٌ، فَإِنْ أُذِنَ
لَكَ وَإِلَّا فَارْجِعْ} --- تین بار اجازت، نہ ملے تو لوٹ جاؤ (بخاری و مسلم)۔
اور: \ar{إِنَّمَا جُعِلَ الِاسْتِئْذَانُ مِنْ أَجْلِ الْبَصَرِ} --- اصل مقصد
\emph{پردے کا تحفظ} ہے (بخاری و مسلم)۔

\medskip
\textbf{(ج) مسلمان کے چھ حقوق (مسلم):} \ar{إِذَا لَقِيتَهُ فَسَلِّمْ عَلَيْهِ}
--- ملو تو سلام کرو \ \textbullet\ \ دعوت دے تو قبول کرو \ \textbullet\ \
نصیحت مانگے تو خیر خواہی کرو \ \textbullet\ \ چھینک کر الحمد للہ کہے تو
\emph{تشمیت} کرو \ \textbullet\ \ بیمار ہو تو عیادت کرو \ \textbullet\ \ فوت ہو
جائے تو جنازے کے ساتھ جاؤ۔

\smallskip
یہی وہ حدیث ہے جس سے ''تشمیت'' کی اصطلاح نکلتی ہے --- اور وہ لازمی سوال میں پوچھی
جا چکی ہے۔
\end{hal}

%---------------------------------------------------------------------
\begin{sawal}{سوال \textenglish{9} \ \textnormal{\small(بہار \textenglish{2022}، سوال \textenglish{5}، \textenglish{6} اور \textenglish{7})}}
احادیث کی روشنی میں حبِّ رسول \saw کی اہمیت بیان کیجیے، نیز اسلام میں شعر و شاعری
کا حکم واضح کریں۔
\end{sawal}

\begin{hal}
\textbf{(الف) حبِّ رسول \saw۔} یہ محبت ایمان کی شرط ہے، اس کا حسن نہیں:

\hadees{لَا يُؤْمِنُ أَحَدُكُمْ حَتَّىٰ أَكُونَ أَحَبَّ إِلَيْهِ مِنْ
وَالِدِهِ وَوَلَدِهِ وَالنَّاسِ أَجْمَعِينَ}{بخاری و مسلم}
\tarjuma{تم میں سے کوئی اس وقت تک مومن نہیں ہو سکتا جب تک میں اسے اپنے باپ،
اپنی اولاد اور تمام لوگوں سے زیادہ محبوب نہ ہو جاؤں۔}

\smallskip
حضرت عمرؓ نے عرض کیا: آپ مجھے جان کے سوا ہر چیز سے زیادہ محبوب ہیں۔ فرمایا:
جب تک جان سے بھی زیادہ محبوب نہ ہو جاؤں۔ عمرؓ نے کہا: اب ایسا ہی ہے۔ فرمایا:
''\ar{الْآنَ يَا عُمَرُ}'' (بخاری)۔ قرآن کا معیار: \ar{النَّبِيُّ أَوْلَىٰ
بِالْمُؤْمِنِينَ مِنْ أَنفُسِهِمْ} (الاحزاب، \textenglish{6})۔

\smallskip
\textbf{تقاضے:} اتباعِ سنت \ \textbullet\ \ کثرتِ درود \ \textbullet\ \ آواز نبی
\saw سے بلند نہ ہو (الحجرات، \textenglish{2}) \ \textbullet\ \ ناموسِ رسالت کا
تحفظ۔ بشارت: \ar{الْمَرْءُ مَعَ مَنْ أَحَبَّ} (بخاری و مسلم)۔

\medskip
\textbf{(ب) شعر و شاعری کا حکم۔} یہ سوال بہار \textenglish{2022} میں دو حصوں
میں آیا --- ایک حدیث کے ترجمے کی صورت میں، ایک مستقل سوال کی صورت میں۔

\hadees{أَصْدَقُ كَلِمَةٍ قَالَهَا الشَّاعِرُ كَلِمَةُ لَبِيدٍ: أَلَا كُلُّ
شَيْءٍ مَا خَلَا اللهَ بَاطِلٌ}{بخاری و مسلم --- حضرت ابو ہریرہؓ \ \ [بہار
\textenglish{2022}]}
\tarjuma{سب سے سچی بات جو کسی شاعر نے کہی، وہ لبید کا یہ مصرع ہے: خبردار! اللہ
کے سوا ہر چیز باطل (فنا ہونے والی) ہے۔}

\smallskip
\textbf{اصل اصول:} شعر بذاتِ خود نہ حلال ہے نہ حرام --- \emph{اس کے مضمون} پر
حکم لگتا ہے۔ نبی کریم \saw نے ایک مشرک شاعر کے سچے مصرعے کو ''سب سے سچی بات''
فرمایا۔

\begin{itemize}\setlength{\itemsep}{1pt}
  \item \textbf{پسندیدہ:} \ar{إِنَّ مِنَ الشِّعْرِ لَحِكْمَةً} --- بعض شعر
        حکمت ہوتے ہیں (بخاری)۔ حسان بن ثابتؓ دربارِ نبوت کے شاعر تھے۔
  \item \textbf{ناپسندیدہ:} \ar{لَأَنْ يَمْتَلِئَ جَوْفُ أَحَدِكُمْ قَيْحًا
        خَيْرٌ لَهُ مِنْ أَنْ يَمْتَلِئَ شِعْرًا} --- جب شعر ہی مقصدِ زندگی بن
        کر قرآن سے غافل کر دے (بخاری و مسلم)۔
  \item \textbf{قرآن کا فیصلہ:} \ar{وَالشُّعَرَاءُ يَتَّبِعُهُمُ الْغَاوُونَ}
        (الشعراء، \textenglish{224}) --- مگر استثنا: \ar{إِلَّا الَّذِينَ
        آمَنُوا وَعَمِلُوا الصَّالِحَاتِ} (\textenglish{227})۔ \textbf{یہ
        استثنا لکھنا ضروری ہے}۔
\end{itemize}

\medskip
\textbf{خلاصہ۔} جو شعر توحید، اخلاق، جہاد یا حق کی حمایت میں ہو --- مستحسن؛ جو
جھوٹ، فحش، شراب، عریانی یا کسی کی ہجو پر مبنی ہو --- حرام؛ اور جو بے ضرر ہو مگر
فرائض سے غافل کر دے --- ناپسندیدہ۔
\end{hal}
